% !TEX root = master.tex
\chapter{Funktionale Umsetzung: Funktion und Skript}
\label{ch:a-funktionen}

Dieses Kapitel ordnet jede geforderte Funktionalit\"at dem Skript bzw.\ der Datei zu, in der sie
umgesetzt ist. Damit ist nachvollziehbar, wo eine bestimmte Funktion zu finden ist
(vgl.\ Abgabevorgabe zur Nachvollziehbarkeit). Tabelle~\ref{tab:a-funktionen} gibt die
vollst\"andige Zuordnung an.

\begingroup
\small
\setlength{\tabcolsep}{5pt}
\renewcommand{\arraystretch}{1.3}
\begin{longtable}{@{}L{4.2cm} L{4.6cm} L{4.7cm}@{}}
\caption{Zuordnung Funktionalit\"at zu Skript (Projekt~A)}
\label{tab:a-funktionen}\\
\toprule
\textbf{Funktionalit\"at} & \textbf{Skript / Datei} & \textbf{Kurzbeschreibung} \\
\midrule
\endfirsthead
\multicolumn{3}{@{}l}{\small\itshape Tabelle \thetable{} -- Fortsetzung}\\
\midrule
\textbf{Funktionalit\"at} & \textbf{Skript / Datei} & \textbf{Kurzbeschreibung} \\
\midrule
\endhead
\midrule
\multicolumn{3}{r@{}}{\footnotesize\itshape Fortsetzung n\"achste Seite}\\
\endfoot
\bottomrule
\endlastfoot

\multicolumn{3}{@{}l}{\textbf{Account, Session und Rechteverwaltung}}\\*
\addlinespace[0.4ex]
Sessionverwaltung & \texttt{init.php} & Session-Start und -Wiederherstellung, CSRF-Token \\
Login / Logout & \texttt{auth.php}, \texttt{admin\_logout.php} & Anmeldung, Rollenweiterleitung, Abmeldung \\
Registrierung & \texttt{auth.php} & Serverseitige Pr\"ufung neuer Konten \\
Automatisches Logout & \texttt{init.php}, \texttt{functions.js} & Timeout-Pr\"ufung (Server) und Warndialog (Client) \\
Benutzerverwaltung (CRUD) & \texttt{admin.php} & Anlegen, \"Andern, Rolle, L\"oschen \\
Rechteverwaltung (Rollen) & \texttt{init.php}, \texttt{admin.php} & Zugriffsschutz \"uber \texttt{is\_admin} \\
Self-Service-Profil & \texttt{user.php} & E-Mail/Passwort \"andern, Account l\"oschen \\
Zentrale Konfiguration & \texttt{config.php} & DB-Zugangsdaten, per \texttt{require\_once} \\
\addlinespace[0.8ex]

\multicolumn{3}{@{}l}{\textbf{Sicherheit}}\\*
\addlinespace[0.4ex]
CSRF-Schutz & \texttt{init.php} (+ Pr\"ufung in \texttt{admin.php}, \texttt{user.php}, \texttt{auth.php}) & Token-Erzeugung und -Pr\"ufung je Formular \\
Passwort-Hashing & \texttt{auth.php} & bcrypt via \texttt{password\_hash}/\texttt{verify} \\
Brute-Force-Schutz & \texttt{auth.php}, Tabelle \texttt{login\_logs} & Sperre nach mehreren Fehlversuchen \\
Rate-Limiting (Registrierung) & \texttt{auth.php}, Tabelle \texttt{register\_logs} & Begrenzung neuer Konten pro IP \\
Prepared Statements & \texttt{auth.php}, \texttt{admin.php}, \texttt{user.php} & Schutz vor SQL-Injection \\
Eingabevalidierung (RegExp) & \texttt{functions.js}, \texttt{check.php}, \texttt{auth.php} & Client- und serverseitig, regul\"are Ausdr\"ucke \\
\addlinespace[0.8ex]

\multicolumn{3}{@{}l}{\textbf{Oberfl\"ache, Bibliotheken und AJAX}}\\*
\addlinespace[0.4ex]
jQuery / jQuery~UI & \texttt{functions.js} & Timeout-Dialog, DOM-Interaktion \\
Bootstrap-Einbindung & \texttt{template/}, \texttt{style/} & Layout- und Komponentenraster \\
AJAX-Nachladen & \texttt{api-calls/check.php}, \texttt{keep\_alive.php} & Live-Pr\"ufung, Session-Verl\"angerung \\
Meldungs-/Best\"atigungsdialoge & \texttt{functions.js}, \texttt{template/footer.php} & \texttt{confirm()} und jQuery-UI-Dialog \\
Dynamische, responsive Navigation & \texttt{template/navbar.php}, \texttt{navbar\_mobile.php} & rollenabh\"angig, \texttt{is\_mobile()}, \texttt{get\_url()} \\
\addlinespace[0.8ex]

\multicolumn{3}{@{}l}{\textbf{Fachmodule und weitere Seiten}}\\*
\addlinespace[0.4ex]
JSON-Datenexport & \texttt{metadata\_stripping.php} & EXIF-Daten als JSON-Download \\
Datei-Upload & \texttt{metadata\_stripping.php} & Bild-Upload und -Bereinigung \\
Externer API-Aufruf & \texttt{api-calls/check.php} & Leak-Abgleich (HaveIBeenPwned) \\
Browser-Fingerprinting & \texttt{fingerprinting/info.php} & Canvas, WebGL, Zeitzone u.\,a. \\
Rechtsseiten & \texttt{impressum.php}, \texttt{datenschutz.php}, \texttt{agb.php} & \"uber den Footer verlinkt \\
\end{longtable}
\endgroup

\section{Sicherheitsmechanismen im Detail}
Die sicherheitsrelevanten Funktionen sind bewusst zentral umgesetzt:

\textbf{CSRF-Schutz.} F\"ur jede Session wird in \texttt{init.php} ein Zufalls-Token erzeugt. Jedes
schreibende Formular f\"uhrt das Token mit; serverseitig wird es zeitkonstant gepr\"uft, bevor eine
Aktion ausgef\"uhrt wird~\autocite{owasp-csrf}.

\begin{lstlisting}[language=PHP,caption={Zeitkonstante CSRF-Pr\"ufung (Auszug \texttt{admin.php})},label={lst:a-csrf}]
if (!hash_equals($_SESSION['csrf_token'], $_POST['csrf_token'])) {
    die("CSRF Token ungültig!");
}
\end{lstlisting}

\textbf{Passwort-Hashing.} Passw\"orter werden nie im Klartext gespeichert, sondern mit dem
bcrypt-Verfahren \"uber \texttt{password\_hash()} gehasht und beim Login mit
\texttt{password\_verify()} gepr\"uft~\autocite{php-passwordhash}.

\textbf{Prepared Statements.} S\"amtliche Datenbankabfragen verwenden parametrisierte Abfragen,
sodass Benutzereingaben nicht als SQL interpretiert werden k\"onnen.

\textbf{Brute-Force- und Rate-Limiting.} Fehlgeschlagene Logins werden in \texttt{login\_logs}
protokolliert und f\"uhren nach mehreren Versuchen zu einer tempor\"aren Sperre; Registrierungen
werden \"uber \texttt{register\_logs} je IP begrenzt.

\textbf{Datenschutzfreundlicher Leak-Abgleich.} Beim Passwort-Check wird nur ein Pr\"afix des
SHA-1-Hashes an die externe API gesendet (k-Anonymity); der Abgleich der vollst\"andigen Hashes
erfolgt lokal~\autocite{hibp-range}.
